% !TeX root = ../main.tex
\section{Pore fabric and shape-changing distention}
\label{sec:pore-fabric}

The spherical specialization of
\cref{sec:kinematics,sec:work-equivalence} lets the distention change pore
volume and rotate the intermediate frame, but not change pore shape. The
fabric literature summarized in \cref{sec:introduction} shows that pore
shape and orientation, not porosity alone, set the effective poroelastic
properties, and that reactions create and orient that structure. This
section introduces a symmetric positive definite distention and derives
its reference-state response. At finite deformation we restrict energy--work
equivalence to states in which distention and mineral stress commute.
This restriction permits a finite pure-shear benchmark with oblique fabric;
it does not establish a general noncoaxial extension. We state at
each stage which assumptions are new modeling choices rather than
consequences of the spherical model.

\subsection{Kinematics of shape-changing distention}
\label{sec:fabric-kinematics}

Write the distention gradient in right polar form as a proper rotation
times a symmetric positive definite stretch:
\begin{equation}
 \mathbf A=\mathbf R_A\mathbf G^{1/2},\qquad
 \mathbf R_A^T\mathbf R_A=\mathbf I,\qquad \det\mathbf R_A=1,\qquad
 \mathbf G=\mathbf G^T,\qquad \mathbf G\ \text{positive definite}.
 \label{eq:fabric-distention-polar}
\end{equation}
Here \(\mathbf G=\mathbf A^T\mathbf A\) is the right Cauchy--Green tensor
of the distention, \(\mathbf G^{1/2}\) its right stretch, and
\(\mathbf R_A\) the distention rotation. The distention volume ratio is
\begin{equation}
 a=\det\mathbf A=\det\mathbf G^{1/2}=(\det\mathbf G)^{1/2}>0.
 \label{eq:fabric-volume-ratio}
\end{equation}
The multiplicative decomposition and its inverse are
\begin{equation}
 \mathbf F=\mathbf A\bar{\mathbf F}
        =\mathbf R_A\mathbf G^{1/2}\bar{\mathbf F},\qquad
 J=a\bar J,\qquad
 \bar{\mathbf F}=\mathbf G^{-1/2}\mathbf R_A^T\mathbf F.
 \label{eq:fabric-multiplicative}
\end{equation}
The mineral right Cauchy--Green tensor follows:
\begin{equation}
 \bar{\mathbf C}=\bar{\mathbf F}^T\bar{\mathbf F}
   =\mathbf F^T\mathbf R_A\mathbf G^{-1}\mathbf R_A^T\mathbf F.
 \label{eq:fabric-mineral-metric}
\end{equation}

The volume-only model is the specialization
\(\mathbf G=a^{2/3}\mathbf I\). Then
\(\mathbf G^{1/2}=a^{1/3}\mathbf I\),
\(\mathbf A=a^{1/3}\mathbf R_A\),
\(\bar{\mathbf F}=a^{-1/3}\mathbf R_A^T\mathbf F\), and
\(\bar{\mathbf C}=a^{-2/3}\mathbf C\), so
\eqref{eq:fabric-multiplicative} and
\eqref{eq:fabric-mineral-metric} reduce exactly to
\eqref{eq:spherical-distention} and
\eqref{eq:spherical-mineral-metric}. In that limit the identity stretch
commutes with \(\mathbf R_A\), so the internal rotation cancels. For a
general \(\mathbf G\) it does not: \(\bar{\mathbf C}\) retains
\(\mathbf R_A\) through \(\mathbf R_A\mathbf G^{-1}\mathbf R_A^T\), and
the relative orientation of the fabric and mineral frames is a physical
datum. For the reference linearization and the finite pure-stretch paths below,
we express the fabric and mineral stiffness in their common reference
axes and set \(\mathbf R_A=\mathbf I\) in
\eqref{eq:fabric-multiplicative}--\eqref{eq:fabric-mineral-metric}. Under a superposed rigid rotation \(\mathbf Q\), both \(\mathbf F\) and
\(\mathbf R_A\) transform on the left by \(\mathbf Q\); setting
\(\mathbf R_A=\mathbf I\) specifies the unrotated calculation, not a
spatially fixed fabric under rotation of the body. The prescribed relative
orientation is a modeling choice; in the spherical specialization it is immaterial
because \(\mathbf R_A\) cancels there. Fixing the rotation has three
consequences. The distention is symmetric in the material frame, so
\(\mathbf G\) carries both the shape and the orientation of the
distention. The orientation of the pore fabric relative to the mineral
therefore enters as prescribed material data rather than as a quantity
the deformation determines. The construction therefore carries no rotational
fabric variable, it represents no rotation of the pore fabric relative to
the mineral matrix, and the elastic equilibrium of
\cref{sec:fabric-equilibrium} dissipates no work in such a rotation. The
drained compliance relation and the Biot tensor derived below are
properties of this fixed-orientation construction, in which the pore
fabric sets the anisotropy of the pressure coupling at a fixed
orientation and only its shape changes. A relative rotation at the
fabric--mineral interface, resisted by friction or by viscous slip, would
require a constitutive relation for \(\mathbf R_A\), the couple conjugate
to its rate, and the work that the resistance dissipates. The fixed
rotation is the limiting case in which the resisting mechanism prevents
relative rotation, and we leave the dissipative case to future work.

\subsection{Mass, volume, and the fabric tensor}
\label{sec:fabric-mass}

The distention volume ratio is \(a=\det\mathbf G^{1/2}\). Separate
\(\mathbf G\) into its volumetric and isochoric parts. The pore-fabric
orientation tensor is the unimodular part
\begin{equation}
 \mathbf H=(\det\mathbf G)^{-1/3}\mathbf G
          =a^{-2/3}\mathbf G,\qquad
 \mathbf H=\mathbf H^T,\qquad \det\mathbf H=1,\qquad
 \mathbf G=a^{2/3}\mathbf H.
 \label{eq:fabric-tensor}
\end{equation}
The scalar \(a\) sets the pore volume change; the unit-determinant tensor
\(\mathbf H\) sets the shape and orientation of the pore space at fixed
pore volume. Its eigenvalues are the squared principal shape ratios of the
distention, and its eigenvectors are the fabric axes. This is the same role the
fabric tensor plays in the elasticity-tensor literature
\citep{cowin1985fabric,turnercowin1987,moesencardosocowin2012}, here
attached to the distention rather than to the stiffness directly.

Conservation of solid mass is unchanged in form because it involves only
the determinant of \(\mathbf A\):
\begin{equation}
 J\rho_s=\phi_{s0}\bar\rho_{s0},\qquad
 \bar J=\frac{\bar\rho_{s0}}{\bar\rho_s},\qquad
 \phi_s=\frac{\phi_{s0}\bar J}{J},\qquad
 \phi_s a=\phi_{s0}.
 \label{eq:fabric-solid-mass}
\end{equation}
Thus \(\phi_s=\phi_{s0}/a\), the current solid volume fraction depends on
the pore volume through the determinant of \(\mathbf G\) alone, and the
fabric \(\mathbf H\) carries no mass. Its effect is purely energetic: it
enters the distention energy below and thereby the stress and the Biot
tensor.

\subsection{Distention energy and its work conjugate}
\label{sec:fabric-energy}

For commuting distention and mineral stress, the energy per reference
mixture volume takes the form
\begin{equation}
 W_s(\mathbf F,\mathbf G)=W_{\mathrm{dis}}(\mathbf G)
   +\phi_{s0}\bar W_s(\bar{\mathbf F}),\qquad
 \bar{\mathbf F}=\mathbf G^{-1/2}\mathbf R_A^T\mathbf F.
 \label{eq:fabric-equivalent-energy}
\end{equation}
The distention energy \(W_{\mathrm{dis}}(\mathbf G)\) now depends on the full symmetric
tensor \(\mathbf G\), replacing the volume-only energy \(W_A(a)\). The
subscript \(\mathrm{dis}\) denotes distention. Its
work conjugate is the symmetric distention stress per reference mixture
volume in the intermediate frame,
\begin{equation}
 \bar{\mathbf S}_{\mathrm{dis}}=2\pder{W_{\mathrm{dis}}}{\mathbf G},\qquad
 \delta W_{\mathrm{dis}}=\tfrac12\bar{\mathbf S}_{\mathrm{dis}}:\delta\mathbf G.
 \label{eq:fabric-distention-stress}
\end{equation}
Equivalently, with the distention logarithmic strain
\(\mathbf E_{\mathrm{dis}}=\ln\mathbf G^{1/2}=\tfrac12\ln\mathbf G\) and
\(\mathbf S_{\mathrm{dis}}=\partial W_{\mathrm{dis}}/\partial\mathbf E_{\mathrm{dis}}\),
\(\delta W_{\mathrm{dis}}=\mathbf S_{\mathrm{dis}}:\delta\mathbf E_{\mathrm{dis}}\); the two stresses are
related by the self-adjoint matrix-logarithm derivative of
\cref{sec:logarithmic-derivative}. The scalar driving term
\(\partial W_A/\partial\ln a=\tfrac13\tr\mathbf\tau'\) of
\eqref{eq:energy-returned-pressure-balance} --- equivalently, the trace of
\eqref{eq:constitutive-kirchhoff-phase-stress} --- is replaced by this
tensor:
its spherical part carries the volume response and its deviatoric part the
shape response.

Carry the variation of
\eqref{eq:fabric-multiplicative} into the coupled work
\eqref{eq:coupled-energy-work}:
\begin{align}
 \delta\mathbf F\mathbf F^{-1}
 &=\delta\mathbf R_A\mathbf R_A^T
   +\mathbf R_A\left[\delta(\mathbf G^{1/2})\mathbf G^{-1/2}\right]
        \mathbf R_A^T
 \nonumber\\
 &\quad+\mathbf R_A\mathbf G^{1/2}
        (\delta\bar{\mathbf F}\bar{\mathbf F}^{-1})
        \mathbf G^{-1/2}\mathbf R_A^T.
 \label{eq:fabric-virtual-deformation}
\end{align}
Contracting with the symmetric effective stress
\(\mathbf\tau'\), and writing
\(\widetilde{\mathbf\tau}=\mathbf R_A^T\mathbf\tau'\mathbf R_A\), removes
the skew rotation increment and, after the phase balance
\eqref{eq:constitutive-kirchhoff-phase-stress} cancels the pressure term,
gives
\begin{align}
 \delta W_s
 &=\widetilde{\mathbf\tau}:
   \left[\delta(\mathbf G^{1/2})\mathbf G^{-1/2}\right]
 \nonumber\\
 &\quad+\phi_{s0}\mathbf G^{1/2}
        \widehat{\mathbf\tau}_s\mathbf G^{-1/2}:
        (\delta\bar{\mathbf F}\bar{\mathbf F}^{-1}).
 \label{eq:fabric-phase-work}
\end{align}
The tilde denotes the effective stress rotated into the true mineral
frame while retaining its reference mixture-volume normalization.
The hatted mineral stress \(\widehat{\mathbf\tau}_s\) is in the same frame
and is normalized per reference mineral volume.
In the spherical limit \(\mathbf G^{1/2}=a^{1/3}\mathbf I\) the first term
becomes \(\tfrac13\tr\mathbf\tau'\,\delta\ln a\) and the second
\(\phi_{s0}\widehat{\mathbf\tau}_s:(\delta\bar{\mathbf F}\bar{\mathbf F}^{-1})\),
recovering \eqref{eq:distention-mineral-energy-work} exactly. The mineral term equals the variation of the prescribed mineral energy
only under the additional restriction
\begin{equation}
 \mathbf G^{1/2}\widehat{\mathbf\tau}_s
 =\widehat{\mathbf\tau}_s\mathbf G^{1/2}.
 \label{eq:fabric-commuting-work}
\end{equation}
Then the similarity transform in \eqref{eq:fabric-phase-work} drops out,
and the remaining work defines the distention contribution in
\eqref{eq:fabric-equivalent-energy}. We impose this restriction for the
finite fabric benchmark. It holds for an isotropic mineral on the coaxial
pure-stretch paths used below. For general noncoaxial states the displayed
energy cannot be inferred from the stated spatial phase balance: the
transformed mineral work would require additional constitutive mechanics.
At the unstressed reference state this discrepancy enters beyond the
quadratic energy, so the reference linearization below remains available
for arbitrary small symmetric strains.

\subsection{Equilibrium and material symmetry}
\label{sec:fabric-equilibrium}

Within the commuting finite specialization, or its reference quadratic
linearization, the retained fabric variables minimize the potential at
fixed deformation and pressure:
\begin{equation}
 \pder{}{\mathbf G}
 \left[W_{\mathrm{dis}}(\mathbf G)+\phi_{s0}\bar W_s(\bar{\mathbf F})
        +\phi_{s0}p\bar J\right]=0,\qquad
 \bar J=J(\det\mathbf G)^{-1/2}.
 \label{eq:fabric-equilibrium}
\end{equation}
This is the tensorial generalization of the pressure equilibrium
\eqref{eq:pressure-conjugacy}. The minimization is restricted to
distentions whose logarithmic strain lies in the retained subspace
\(\operatorname{range}\mathbb{D}\) of \cref{sec:fabric-biot}. On the
complement \(\mathbb{D}\) vanishes, so the distention energy is independent
of the complementary modes; those modes are dropped from the minimization by
construction and frozen to the mineral rather than equilibrated. Equation
\eqref{eq:fabric-equilibrium} is therefore read on that subspace, where the
retained quadratic distention potential is strictly convex and the mineral term
contributes a positive-definite stiffness in the same strain, so the stationary
point exists and is unique in the reference quadratic model of
\cref{sec:fabric-biot}. Differentiating along the spherical
submanifold \(\mathbf G=a^{2/3}\mathbf I\), on which only \(a\) varies,
reduces \eqref{eq:fabric-equilibrium} to
\(\partial W_A/\partial\ln a
=\tfrac{\phi_{s0}}3\tr\bar{\mathbf\tau}_s+\phi_{s0}p\bar J\), which is
\eqref{eq:energy-returned-pressure-balance}. The remaining retained
components of \eqref{eq:fabric-equilibrium} are new: in the
volumetric--axial model of \cref{sec:fe-fabric} one shape component
equilibrates the
fabric shape against the mineral stress, with no external conjugate force.
Treating the fabric as an elastically equilibrated, nondissipative
internal variable is a new modeling choice; a viscous or reaction-driven
evolution law for \(\mathbf H\) would lie outside the present elastic
scope.

Objectivity restricts \(W_{\mathrm{dis}}\). Under a superposed spatial rotation
\(\mathbf A\mapsto\mathbf Q\mathbf A\) the tensor
\(\mathbf G=\mathbf A^T\mathbf A\) is invariant, so \(W_{\mathrm{dis}}(\mathbf G)\)
and \eqref{eq:fabric-equivalent-energy} are spatial-objective. The
distention stress \eqref{eq:fabric-distention-stress} is symmetric. The prescribed material fabric axis is a structural datum of the
reference distention law. In the transversely isotropic specialization,
\(W_{\mathrm{dis}}\) may depend on both \(\mathbf G\) and
\(\mathbf m\otimes\mathbf m\), expressed in the common material axes.
Spatial objectivity does not require material isotropy. An energy depending
only on the invariants of \(\mathbf G\) would be materially isotropic and
would not supply a preferred axis at \(\mathbf G=\mathbf I\).
The prescribed axis therefore supplies the symmetry of the retained
volumetric--axial stiffness, while the equilibrated eigenvalues describe
its shape response.

\subsection{The Biot tensor with fabric coupling}
\label{sec:fabric-biot}

The Biot tensor retains the form
\eqref{eq:finite-biot-tensor},
\begin{equation}
 \mathbf B=\mathbf I-\frac{\phi_{s0}}J
   \left.\pder{\bar J}{\mathbf F}\right|_p\mathbf F^T,
 \label{eq:fabric-biot-tensor}
\end{equation}
but \(\bar J=J(\det\mathbf G)^{-1/2}\) now varies with \(\mathbf F\)
both directly through \(J\) and through the equilibrated fabric
\(\mathbf G(\mathbf F,p)\). The reference result is a symmetric tensor whose directional pressure
response carries both mineral anisotropy and fabric anisotropy. The finite
benchmark evaluates this derivative only at states satisfying
\eqref{eq:fabric-commuting-work}. In the spherical model the shape terms were absent because
\(\mathbf G\) had only the volume degree of freedom; here they are present
and are determined by the distention stiffness.

The spherical restriction is rank one because the distention energy there
depends only on the volume. For a fabric distention the retained subspace is
wider, so the restriction must be derived again. In
the reference state \(\mathbf F=\mathbf G=\mathbf I\), linearize
\eqref{eq:fabric-multiplicative} and \eqref{eq:fabric-mineral-metric}. To
first order the logarithmic strains add:
\begin{equation}
 \bar{\mathbf\varepsilon}=\mathbf\varepsilon-\mathbf E_{\mathrm{dis}}.
 \label{eq:fabric-additive-strain}
\end{equation}
With a quadratic reference distention energy
\(W_{\mathrm{dis}}=\tfrac12\mathbf E_{\mathrm{dis}}:\mathbb{D}:\mathbf E_{\mathrm{dis}}\), the
distention stiffness tensor \(\mathbb D\) is symmetric and positive
semidefinite: positive definite on the subspace of distention strains the
model retains and zero on its complement. The minimization in
\eqref{eq:fabric-equilibrium} is therefore over
\(\mathbf E_{\mathrm{dis}}\in\operatorname{range}\mathbb{D}\), not over all
symmetric tensors; the complementary modes are kinematically frozen to the
mineral, which carries them at its own compliance. On that subspace
\(\mathbb{D}+\phi_{s0}\mathbb{C}_s\) is positive definite, which is the condition
for the unique stationary point of \eqref{eq:fabric-equilibrium} used above;
minimizing at fixed \(\mathbf\varepsilon\) gives the drained compliance
\begin{equation}
 (\mathbb{C}^d)^{-1}
 =(\phi_{s0}\mathbb{C}_s)^{-1}+\mathbb{D}^{+}.
 \label{eq:fabric-compliance-restriction}
\end{equation}
with \(\mathbb D^{+}\) the Moore--Penrose inverse: the ordinary inverse on
\(\operatorname{range}\mathbb{D}\), extended by zero on the complement.
Here \(\mathbb{D}\) has major and minor symmetry and the material symmetry
of the fabric. In the spherical case \(W_{\mathrm{dis}}\) depends only on
\(\tr\mathbf E_{\mathrm{dis}}=\ln a\), so \(\mathbb{D}\propto\mathbf I\otimes\mathbf I\)
is rank one and \eqref{eq:fabric-compliance-restriction} reduces to
\eqref{eq:drained-compliance-restriction}, with \(\mathbb{D}^{+}\) the
corresponding rank-one inverse on the volume direction. Retaining a
wider subspace relaxes the additional compliance from a spherical rank-one
tensor to a fabric-symmetric fourth-order tensor supported on that subspace.
The model implemented in \cref{sec:fe-fabric} retains the two-dimensional
axisymmetric subspace spanned by the volumetric direction
\(\mathbf e_1=\mathbf I/\sqrt3\) and the axial (degree-two) direction
\(\mathbf e_2=\sqrt{3/2}\,(\mathbf m\otimes\mathbf m-\mathbf I/3)\),
where \(\mathbf m\) is the unit material fabric axis; its
distention stiffness is the positive definite \(2\times2\) block
\((k_v,k_a,k_c)\) in that basis, of rank two on the six-dimensional space of
symmetric tensors, and the four complementary modes are frozen to the
mineral compliance through the zero eigenvalues of \(\mathbb D^{+}\). Let
\(\mathbf p_1,\mathbf p_2\) be orthonormal vectors spanning the plane normal
to \(\mathbf m\). That reduced \(\mathbb{D}\) is transversely isotropic
about \(\mathbf m\) because \(\mathbf e_1\) and \(\mathbf e_2\) are
invariant under rotations about \(\mathbf m\); and, concretely, it
annihilates the four complementary modes of this basis, the coupled
in-plane pair
\(\mathbf e_3=(\mathbf p_1\otimes\mathbf p_1-\mathbf p_2\otimes\mathbf p_2)/\sqrt2\)
and \(\mathbf e_6=\sqrt2\operatorname{sym}(\mathbf p_1\otimes\mathbf p_2)\)
together with the two axial-shear modes
\(\mathbf e_4=\sqrt2\operatorname{sym}(\mathbf m\otimes\mathbf p_1)\) and
\(\mathbf e_5=\sqrt2\operatorname{sym}(\mathbf m\otimes\mathbf p_2)\). The labels
\(\mathbf e_1,\dots,\mathbf e_6\) denote this fabric-adapted orthonormal
basis of symmetric tensors, not the Mandel component indices of
\eqref{eq:example-mineral-stiffness}. Two stiffnesses
selected independently therefore satisfy the generalized restriction when
their difference is this distention compliance, not only a spherical
contribution. The reference Biot tensor
\(\mathbf B_0=\mathbf I-\mathbb{C}^d:\mathbb{C}_s^{-1}:\mathbf I\) of
\eqref{eq:reference-biot-compatibility} keeps its form with the
generalized \(\mathbb{C}^d\) and becomes a general anisotropic tensor;
its couplings arise from both the mineral and the fabric.

\subsection{Transverse isotropy of the pore fabric}
\label{sec:fabric-transverse}

Let the fabric be transversely isotropic about a unit material axis
\(\mathbf m\), and restrict the distention logarithmic strain to the
axisymmetric subspace \(\operatorname{span}(\mathbf e_1,\mathbf e_2)\) of
\cref{sec:fabric-biot}. Then
\begin{equation}
 \mathbf H=h^{-2}\mathbf m\otimes\mathbf m
   +h(\mathbf I-\mathbf m\otimes\mathbf m),\qquad h>0,
 \label{eq:fabric-transverse-h}
\end{equation}
so \(\det\mathbf H=1\) and
\(\mathbf G=a^{2/3}\left[h^{-2}\mathbf m\otimes\mathbf m
+h(\mathbf I-\mathbf m\otimes\mathbf m)\right]\). The distention has two
scalar degrees of freedom, the volume ratio \(a\) and the fabric shape
parameter \(h\), and \(W_{\mathrm{dis}}=W_{\mathrm{dis}}(a,h)\). With the
strain in that subspace the inverse relation is exact: the distention
logarithmic strain is
\begin{equation}
 \mathbf E_{\mathrm{dis}}=\frac{\ln a}{3}\mathbf I
   +\ln h\left(\tfrac12\mathbf I-\tfrac32\mathbf m\otimes\mathbf m\right),
 \label{eq:fabric-transverse-strain}
\end{equation}
twice whose exponential, \(\exp(2\mathbf E_{\mathrm{dis}})=\mathbf G
=a^{2/3}\mathbf H\), reconstructs \eqref{eq:fabric-transverse-h},
whose unimodular eigenvalues are \(h^{-2},h,h\). The retained model of
\cref{sec:fe-fabric} holds exactly in this subspace, so the reported shape
scalar is exactly \(\ln h\), the logarithm of the unimodular transverse
fabric eigenvalue, and not a projection of \(h\).
The equilibrium \eqref{eq:fabric-equilibrium} splits into a volume balance
for \(a\), which couples to pressure and reduces to the spherical pressure
equation, and a shape balance for \(h\),
\begin{equation}
 \pder{}{h}\left[W_{\mathrm{dis}}+\phi_{s0}\bar W_s\right]=0,
 \label{eq:fabric-shape-equilibrium}
\end{equation}
which has no external conjugate. Because \(\bar J=J/a\) is independent of
\(h\), the fabric shape is set by the balance between the distention
stiffness and the mineral stress, not by the fluid pressure directly.

For an isotropic mineral the drained compliance
\eqref{eq:fabric-compliance-restriction} is then transversely isotropic,
with a distention stiffness \(\mathbb D\) described by the volumetric and
axial-shape fabric moduli. The Biot tensor takes the form
\begin{equation}
 \mathbf B=B_\parallel\,\mathbf m\otimes\mathbf m
   +B_\perp(\mathbf I-\mathbf m\otimes\mathbf m),
 \label{eq:fabric-transverse-biot}
\end{equation}
with distinct axial and transverse coupling coefficients \(B_\parallel\)
and \(B_\perp\). Their difference measures how the axial and in-plane pore
opening resistances differ and vanishes when the fabric is isotropic
(\(h\equiv1\) with \(W_{\mathrm{dis}}\) independent of \(h\)), in which case
\eqref{eq:fabric-transverse-biot} returns the scalar Biot coefficient of
the spherical model. The extension therefore predicts that a
reaction-induced pore fabric aligned along a preferred axis can produce a
directionally asymmetric pressure coupling even when the mineral itself is
isotropic. In the retained volumetric--axial energy used in
\cref{sec:fe-fabric} that
directional coupling is carried by the volume--axial coupling modulus of the
distention stiffness: with that modulus zero, the drained
compliance is transversely isotropic for an isotropic mineral while
\(\mathbf B\) stays spherical.
